% [NEW]: Tom tat tieng Viet - viet moi hoan toan
\begin{center}
\textbf{\large{Tóm tắt}}
\end{center}

\addcontentsline{toc}{chapter}{Tóm tắt}

\begin{small}

% [MODIFIED]: Chinh lai mo ta kien truc tom tat de dong bo voi Chuong 3 va code hien tai; bo cach mo rong DCP khong con phu hop va bo sung vai tro cua tac tu Hieu chinh.
Việc tự động sinh câu hỏi trắc nghiệm từ tài liệu giáo dục là nhu cầu cấp thiết trong bối cảnh số hóa giáo dục, nhưng các giải pháp hiện tại phụ thuộc vào mô hình ngôn ngữ lớn thương mại với chi phí cao và yêu cầu kết nối internet ổn định. Khóa luận này đề xuất kiến trúc đa tác tử, sử dụng các mô hình ngôn ngữ lớn cỡ nhỏ (sLLM - Small Large Language Model) để sinh câu hỏi trắc nghiệm chất lượng cao từ ảnh chụp văn bản. Hệ thống triển khai ba tác tử \textit{sinh câu hỏi} hoạt động song song theo từng cấp độ Bloom, kết hợp với tác tử \textit{Phân loại} để kiểm chứng cấp độ nhận thức, tác tử \textit{Học sinh} để kiểm tra khả năng giải được và tác tử \textit{Hiệu chỉnh} để sửa các phương án trả lời chưa đạt yêu cầu. Khung đánh giá ba tiêu chí gồm khả năng giải được dựa trên ngữ cảnh, chất lượng phương án nhiễu và độ bám sát phân loại nhận thức, kết hợp với phương pháp sử dụng LLM để đánh giá (LLM-as-a-Judge), được xây dựng để đánh giá tự động chất lượng câu hỏi. Ngoài ra, quy trình sinh dữ liệu OCR sử dụng kỹ thuật tạo ảnh tổng hợp bằng trình duyệt nhằm phát triển bộ dữ liệu đánh giá. Kết quả thực nghiệm cho thấy kiến trúc đa tác tử cải thiện chất lượng sLLM rõ rệt, với mức tăng tương đối 93,6\% ở chỉ số chất lượng phương án nhiễu của Phi-4 và chi phí dưới \$0.50 cho mỗi 100 đoạn văn. Xét theo tỷ lệ chấp nhận (Acceptance), cấu hình Phi-4 kết hợp đa tác tử đạt trung bình 0,8557 trên ba cấp độ, trong đó đạt 0,9753 ở Cấp 1 và 0,9072 ở Cấp 2, cạnh tranh sát sao với Gemini 3 Flash ở hai cấp độ này; tuy nhiên, nhóm LLM thương mại vẫn duy trì ưu thế ở Cấp 3. Ứng dụng Openlet được phát triển để minh họa khả năng ứng dụng thực tế của hệ thống.

\vspace*{1cm}
\textbf{Từ khóa}: Mô hình ngôn ngữ lớn cỡ nhỏ, Hệ thống đa tác tử, Câu hỏi trắc nghiệm, Nhận dạng ký tự quang học, Phân loại Bloom, LLM-as-a-Judge

\end{small}
